\documentclass[11pt,a4paper]{article}
\usepackage{fontspec,amsmath,amssymb,amsthm,microtype}
\setmainfont{lmroman10-regular.otf}[BoldFont=lmroman10-bold.otf,ItalicFont=lmroman10-italic.otf,BoldItalicFont=lmroman10-bolditalic.otf]
\usepackage[margin=25mm]{geometry}
\usepackage[unicode,hidelinks]{hyperref}
\hypersetup{pdfauthor={LangevinLean Project}}
\setlength{\emergencystretch}{3em}
\setlength{\parskip}{0.4em}
\setlength{\parindent}{0pt}
\allowdisplaybreaks[2]
\raggedbottom
\title{\texorpdfstring{Fresh Gaussian innovations with retained initial data}{Fresh Gaussian innovations with retained initial data}}
\author{LangevinLean Project}
\date{}
\begin{document}\maketitle
Draft mathematical exposition; these notes may contain mistakes.


The weak consistency proof needs the actual joint law of the current recursive state and its next innovation, including a passive initial weight. These results obtain that law from independence of initialization and an iid noise path.



For iid innovations of law \(\eta\), independent of initialization, the actual recursion has the fresh joint law


\[X_{k+1}=F(X_k,W_k),\qquad
\operatorname{Law}(X_k,W_k)=\operatorname{Law}(X_k)\otimes\eta.\]


Copying a passive initial coordinate through each update retains it exactly:


\[\operatorname{Chain}_{\widetilde F}((a,x),w)_n
=\bigl(a,\operatorname{Chain}_F(x,w)_n\bigr),\qquad
\widetilde F((a,x),z)=(a,F(x,z)).\]

Earlier results used here: \cite[Section~3.33]{LangevinLean2026ULASC}.

\section{\texorpdfstring{Identify the current-state and next-noise law}{Identify the current-state and next-noise law}}


The iid family has the infinite product path law. Independence makes its joint law with the initialization a product. Push this identity through the measurable recursion and use freshness of the canonical product chain to obtain the product of the actual current-state law and the innovation law.



Earlier results used here: \cite[Section~3.20.1]{LangevinLean2026ULASC}; \cite[Section~3.20.4]{LangevinLean2026ULASC}; \cite[Section~3.20.8]{LangevinLean2026ULASC}; \cite[Section~3.33.7]{LangevinLean2026ULASC}.

\section{\texorpdfstring{Retain the initial weight exactly}{Retain the initial weight exactly}}


Run the recursion on a pair whose first coordinate is copied at each step. Induction proves that this coordinate remains the original value while the second coordinate is precisely the ordinary noise-driven chain. This allows a bounded initial weight to be included in the freshness argument.



Earlier results used here: \cite[Section~3.20.1]{LangevinLean2026ULASC}; \cite[Section~3.20.3]{LangevinLean2026ULASC}.

\begin{thebibliography}{ULA-LSI}
\bibitem[ULA-SC]{LangevinLean2026ULASC}
LangevinLean Project. \emph{Unadjusted Langevin under Strong Convexity}. LL-01, 2026. Draft report, content revision c58ab5c70c3a. \url{https://odakyildiz.github.io/LangevinLean/docs/pdf/UnadjustedLangevin/FullProof.pdf}.
\end{thebibliography}

\end{document}
